\documentclass[11pt]{article}
\usepackage[a4paper,margin=23mm]{geometry}
\usepackage[T1]{fontenc}
\usepackage{lmodern,microtype,amsmath,amssymb,bm,booktabs,array,longtable,xurl,hyperref}
\hypersetup{colorlinks=true,linkcolor=blue,urlcolor=blue}
\allowdisplaybreaks\raggedbottom
\newcommand{\Tr}{\operatorname{Tr}}
\newcommand{\dd}{\mathrm d}
\newcommand{\MT}{\mathrm{MT}}
\title{Audit of the remaining LAPW stress terms:\\core Eq. (E3), electrostatics, and GGA\\with a SrVO$_3$ band-completeness test}
\author{Working derivation and implementation record for Kristjan Haule}
\date{7 November 2026}
\begin{document}
\maketitle
\begin{abstract}
This audit finds and corrects three implementation defects: signed lattice-harmonic labels incorrectly used as angular momenta in core and electrostatic contractions; incomplete GGA gradient tensors, including missing mixed-spin terms; and a single-precision conversion of a large electrostatic contribution. Independent polynomial and native-PBE finite-difference tests establish these defects without fitting a material pressure. The corrected local executables are installed. These repairs do \emph{not} resolve the SrVO$_3$ RKMAX=8 eDMFT mismatch. Increasing the number of stored bands changes its analytic stress significantly but moves it farther from the numerical derivative. A separate GGA functional-consistency residual also remains. This note distinguishes verified identities, isolated regression tests, and unresolved material discrepancies.
\end{abstract}
\tableofcontents
\clearpage

\section{Scope, units, and comparison convention}
The source paper is K. Belbase, A. Tr\"oster, and P. Blaha, \emph{Phys. Rev. B} \textbf{104}, 174113 (2021), supplied locally as \nolinkurl{Blahas_stress_tensor.pdf}. This pass concentrates on Appendix E and the density-dependent stress in LAPW0; it does not repeat the two preceding G1/G2 audits.

Let $F=I+\epsilon$ deform the cell, $\Omega$ denote the reference volume, and $\sigma_{ab}=\Omega^{-1}\partial\Gamma/\partial\epsilon_{ab}$ denote stress, where $\Gamma$ is the stationary free-energy functional. Cartesian indices are $a,b$. For an isotropic linear strain $F=(1+t)I$,
\begin{equation}
 \frac{\dd\Gamma}{\dd t}=\Omega\Tr\sigma,\qquad
 \bar\sigma=\frac{\Tr\sigma}{3},\qquad p=-\bar\sigma.
\end{equation}
All numerical results below are \textbf{mean normal stress} $\bar\sigma$ in GPa. A difference quoted here is already divided by three. The conversion used in the tests is $1\ \mathrm{Ry}/a_0^3=14710.5164057$ GPa.

The WIEN intermediate arrays have the assembly convention
\begin{equation}
 \sigma^{\rm total}=S^{\rm DENS}-S^{\rm CORE}-S^{\rm VAL}.
\end{equation}
Consequently a positive stored core tensor is subtracted in the physical stress. This sign must not be applied twice. The density terms below are evaluated on the supplied density; their formulas apply equally to a DFT or an eDMFT density. The remaining eDMFT matrix terms and normalized-projector derivatives are separate.

Tests used isolated copies of the central SrVO$_3$ DFT/eDMFT and FeSe DFT/eDMFT inputs. The converged reference directories were not edited or rerun. The current FeSe input snapshots differ from some earlier archived comparisons, so their fixed-density changes below must not be silently substituted into an older total-stress table.

\section{Core correction: derivation of the angular content of E3}
Let $A$ label a nucleus, $\bm r$ be displacement from that nucleus, $R_A$ its muffin-tin radius, $r=|\bm r|$, and $\bm n=\bm r/r$. The spherical core number density is $\rho_c^A(r)$. Introduce
\begin{equation}
 q_A(r)=4\pi r^2\rho_c^A(r),\qquad
 V_A(\bm r)=\sum_{lm}V^A_{lm}(r)Y_{lm}(\bm n).
\end{equation}
The complex spherical harmonics obey $\int |Y_{lm}|^2\dd\omega=1$; $\dd\omega$ is solid angle. All potentials in the equations are Rydbergs. A prime on a radial coefficient denotes $\dd/\dd r$.

The core virial correction represented by Appendix E is
\begin{equation}\label{eq:core}
 \sigma^{c}_{ab}=-\frac1{2\Omega}\sum_A\int_{r<R_A}\rho_c^A(r)
       (r_b\partial_a+r_a\partial_b)V_A(\bm r)\,\dd^3r.
\end{equation}
Using the dimensionless tangential derivative
\begin{equation}
 D_a=(\delta_{aj}-n_an_j)\frac{\partial}{\partial n_j},\qquad
 \partial_a=n_a\partial_r+r^{-1}D_a,
\end{equation}
we obtain the stored positive core contribution
\begin{equation}\label{eq:radialcore}
 \Omega S^{\rm CORE}_{ab}=
 \sum_{A,lm}\int_0^{R_A}q_A(r)
 \left[r(V^A_{lm})'(r)\,\mathcal A^{lm}_{ab}
       +V^A_{lm}(r)\,\mathcal B^{lm}_{ab}\right]\dd r,
\end{equation}
where
\begin{align}
 \mathcal A^{lm}_{ab}&=\frac1{4\pi}\int n_an_bY_{lm}\,\dd\omega,\\
 \mathcal B^{lm}_{ab}&=\frac1{8\pi}\int(n_bD_a+n_aD_b)Y_{lm}\,\dd\omega.
\end{align}
This is the radial/angular decomposition underlying E3. Both angular averages carry the same $1/(4\pi)$ normalization after symmetrization.

\subsection{An independent angular identity}
On the unit sphere, $D_an_b=\delta_{ab}-n_an_b$ and the tangential vector field associated with $D_a$ has divergence $-2n_a$. Integration by parts therefore gives
\begin{equation}
 \int n_bD_aY\,\dd\omega=
 \int(3n_an_b-\delta_{ab})Y\,\dd\omega.
\end{equation}
Hence
\begin{equation}\label{eq:angularidentity}
 \boxed{\mathcal B^{lm}_{ab}=3\mathcal A^{lm}_{ab}
 -\delta_{ab}\langle Y_{lm}\rangle_\omega.}
\end{equation}
Only $l=0$ and $l=2$ survive. For $l=0$, $\mathcal B=0$; for $l=2$, $\mathcal B=3\mathcal A$ and $\Tr\mathcal A=0$. Writing the physical spherical potential as $\bar V=V_{00}/\sqrt{4\pi}$ yields
\begin{equation}
 \Omega S^{{\rm CORE},0}_{ab}=
 \frac{\delta_{ab}}3\sum_A\int_0^{R_A}q_A(r)r\bar V'_A(r)\,\dd r.
\end{equation}
This provides a scalar check independent of the Gaunt implementation.

\subsection{Confirmed core defect and its test}
WIEN uses a negative stored $l$ to select a sine lattice harmonic. It is \emph{not} a negative angular momentum. After the real harmonic has been expanded into complex coefficients, Gaunt triangle rules and derivative coefficients require $|l|$. Two loops in \nolinkurl{SRC_lcore/core_corr_stress.f} still used signed $l$ and silently discarded valid sine quadrupoles. They now use $|l|$, with the sine phase retained in the complex coefficients.

An independent Cartesian fixture takes $q=r^2$, $R=1$, and $V=\bm r^TQ\bm r$. Equation~\eqref{eq:core} gives the stored result $2Q/15$. It tests the spherical component, all five real quadrupoles, and vanishing $l=1,3,4$ contributions. The former code returned zero for two shear components whose correct values were $0.07283656204$. The corrected maximum error over the fixture is $1.51\times10^{-8}$, set by radial quadrature. These $l=2$ core terms are traceless and cannot explain a hydrostatic SrVO$_3$ discrepancy.

A separate check interpolated the smooth part of each actual spherical potential after removing its nuclear $-2Z/r$ pole and differentiated the pole analytically. Its change relative to the stored core mean was about $0.00290$ GPa in SrVO$_3$ and $0.00334$ GPa in FeSe, nearly identical for their DFT and eDMFT densities. No large core radial-derivative error was demonstrated. These checks do not prove away every frozen-core approximation or justify adding an isolated origin or boundary term without its electrostatic cancellation partner.

\section{Electrostatic stress and precision}
The same signed-$l$ defect occurred in the LAPW0 muffin-tin electrostatic contractions. Here both density and potential can be nonspherical. Products of two sine harmonics can have a nonzero scalar part, so this defect can affect the trace as well as shear. The relevant indices in \nolinkurl{SRC_lapw0/finl_elect_str.F} now use $|l|$ after conversion to complex harmonics.

A fixture compares the actual Fortran contraction with independent Cartesian angular quadrature of
\begin{equation}
 -\frac12\int_{r<1}\rho(\bm r)V(\bm r)n_an_b\,\dd^3r.
\end{equation}
It exercises a quadrupole in $V$, a quadrupole in $\rho$, and their products, including sine components on either side. The original maximum error was $0.0915291$ in fixture units. The corrected maximum error is $1.68\times10^{-8}$.

The final assembly also called \texttt{REAL(col\_pot\_1)} although that array was already double precision. For a real argument, \texttt{REAL} without a kind converts to default single precision. This is different from extracting the real part of a double-complex argument. The conversion is removed so that large density, core, and valence terms remain in double precision until their cancellation.

\subsection{A combined hydrostatic identity}
Under a normalized uniform scaling, $\rho_s(s\bm r)=s^{-3}\rho(\bm r)$, with nuclei and cell scaled together, the full Coulomb energy obeys $E_C(s)=E_C(1)/s$. Thus
\begin{equation}\label{eq:coultrace}
 \boxed{\Omega\Tr\sigma_C=-E_C.}
\end{equation}
This checks the \emph{sum} of the sphere, interstitial, and nuclear contributions. In LAPW0's KXC output, $E_C=\texttt{DEN}-\texttt{INTRV}-\texttt{EXC}$. Here the quantity printed as \texttt{INTRV} is the kinetic-energy correction $-\int\rho V_{\rm eff}$, despite the abbreviated output label.

The residual $(\Tr\sigma_C+E_C/\Omega)/3$ is:
\begin{center}
\begin{tabular}{lrr}\toprule
Density snapshot & Before (GPa) & Corrected (GPa)\\\midrule
SrVO$_3$ DFT & $-0.001983$ & $-0.001478$\\
SrVO$_3$ eDMFT & $-0.003825$ & $-0.001478$\\
FeSe DFT & $-0.065402$ & $-0.000452$\\
FeSe eDMFT & $-0.057307$ & $-0.000452$\\\bottomrule
\end{tabular}
\end{center}
The corrected residual is almost density-independent between DFT and eDMFT for each material. This strongly constrains a hydrostatic electrostatic explanation of the large eDMFT mismatch. It is not a complete six-component finite-strain validation.

\section{The full spin-resolved GGA metric derivative}
Let $n_u,n_d$ be spin densities, $\bm g_s=\nabla n_s$, and let $f(n_u,n_d,\sigma_{uu},\sigma_{dd},\sigma_{ud})$ be the XC energy density in Ry/$a_0^3$, with
\begin{equation}
 \sigma_{uu}=|\bm g_u|^2,\quad \sigma_{dd}=|\bm g_d|^2,\quad
 \sigma_{ud}=\bm g_u\cdot\bm g_d.
\end{equation}
Define $a_u=\partial f/\partial\sigma_{uu}$, $a_d=\partial f/\partial\sigma_{dd}$, and $a_{ud}=\partial f/\partial\sigma_{ud}$. Gradients transform as $\bm g\mapsto F^{-T}\bm g$. Their explicit metric derivative is
\begin{equation}\label{eq:gga}
 \boxed{Q_{ab}=2a_u g_{u,a}g_{u,b}+2a_d g_{d,a}g_{d,b}
       +a_{ud}(g_{u,a}g_{d,b}+g_{d,a}g_{u,b}).}
\end{equation}
After combining density response with the local potential, the usual XC stress is
\begin{equation}\label{eq:xctotal}
 \Omega\sigma^{xc}_{ab}=\delta_{ab}\int\left[f-\sum_s n_sv_s\right]\dd^3r
 -\int Q_{ab}\,\dd^3r,
 \qquad v_s=\frac{\delta E_{xc}}{\delta n_s}.
\end{equation}
The second integral has a minus sign. It contains no extra density or occupation factor. In an unpolarized representation with total $n=n_u+n_d$, $Q_{ab}=2(\partial f/\partial|\nabla n|^2)g_ag_b$.

\subsection{What the old code omitted}
The muffin-tin implementation contracted same-harmonic radial density derivatives. This is not the full product of Cartesian gradients: it omits cross-harmonic products and tangential derivatives, and it projects a nonlinear gradient coefficient before multiplication. The magnetic interstitial expression used the total gradient in place of the separate spin gradients and omitted the mixed-spin derivative.

The corrected code evaluates Eq.~\eqref{eq:gga} at the existing XC angular points and integrates all nine tensor components before equivalent-atom rotation. GRANS supplies components in $(\bm e_r,\bm e_\theta,\bm e_\phi)$; these are first converted to fixed Cartesian axes. The radial integrand is $r^2\int Q_{ab}\dd\omega$. Interstitial evaluation uses the same algebra.

The existing VXCLM2 interface returns \texttt{sig}=$2\,\partial f_{\rm Ha}/(\mathrm{JSPIN}\,\partial\sigma)$. Thus $a_{\rm Ry}=\mathrm{JSPIN}\,\texttt{sig}$ for the spin-resolved interface. For JSPIN=1 the independent invariant is the total-density gradient squared. XCPOT3 subsequently multiplies its interstitial buffers by $2n_s$; the new buffer contains $Q/[2(n_u+n_d)]$ in each spin slot so that the existing sum returns exactly $Q$.

The helper \nolinkurl{gga_metric_tensor.f} documents these conventions. A diagnostic \nolinkurl{WIEN_STRESS_GGA_MODE=LEGACY} retains the former GGA contraction for comparisons. The normal default is the full tensor. Stress with the alternate \texttt{.xcpot1q} route now fails explicitly, because that route does not supply the required gradient data.

\subsection{Tests and remaining XC limitation}
The helper passes all six symmetric strain derivatives for 60 arbitrary quadratic spin-gradient energies, with maximum error $2.56\times10^{-10}$. Rotation covariance and equivalent spin representations agree to about $7.1\times10^{-15}$. A second test differentiates the actual WIEN PBEX/PBEC energy for 40 spin-polarized density/gradient samples; it agrees with Eq.~\eqref{eq:gga} to $1.15\times10^{-11}$.

A whole-cell XC test scales the density, cell, muffin-tin radii, and radial grids together. Stored radial shell coefficients scale as $s^{-1}$ and interstitial density coefficients as $s^{-3}$. A five-point derivative using $t=\pm0.002,\pm0.004$ is compared with Eq.~\eqref{eq:xctotal}. The residual in mean stress remains:
\begin{center}
\begin{tabular}{lr}\toprule
Density snapshot & Analytic minus scaling derivative (GPa)\\\midrule
SrVO$_3$ DFT (LDA) & $+0.01244$\\
SrVO$_3$ eDMFT (LDA) & $+0.01244$\\
FeSe DFT (GGA) & $+0.37945$\\
FeSe eDMFT (GGA) & $+0.36776$\\\bottomrule
\end{tabular}
\end{center}
\textbf{This full functional-consistency test does not yet pass.} The native point-energy metric test does pass, so replacing the exact tensor with the former incomplete expression would not be justified. Remaining candidates include radial differentiation, angular projection, piecewise-density boundary consistency, and the Euler-potential integration by parts used in Eq.~\eqref{eq:xctotal}. The existing angular derivative treatment at polar points is also a limitation for general harmonics; it was not changed here. No empirical pressure offset has been introduced.

\section{Effect of the installed fixes on the central density contribution}
These are changes at the same density, not new self-consistent total-stress comparisons:
\begin{center}
\begin{tabular}{lrrr}\toprule
Snapshot & Coulomb/precision & Full GGA increment & Total shift\\
 & (GPa) & (GPa) & (GPa)\\\midrule
SrVO$_3$ DFT & $0.000505$ & $0$ & $0.000505$\\
SrVO$_3$ eDMFT & $0.002347$ & $0$ & $0.002347$\\
FeSe DFT & $0.064950$ & $0.184193$ & $0.249143$\\
FeSe eDMFT & $0.056855$ & $0.179044$ & $0.235899$\\\bottomrule
\end{tabular}
\end{center}
The core tensor is unchanged on these four actual snapshots, although the independent sine-quadrupole fixture changes as required. At one OpenMP thread, ordinary $V_{\rm sp}$, $V_{\rm ns}$, and core-density files remain byte-for-byte identical. Two-thread stress agrees at printed precision; nonspherical-potential differences are at most $10^{-13}$ Ry. MPI tests with four ranks for SrVO$_3$ and two for FeSe reproduce the serial tensor at printed precision.

\section{\texorpdfstring{SrVO$_3$}{SrVO3}: stored bands are a separate convergence parameter}
The RKMAX=8 central calculation has 75 irreducible $k$ points and LAPW matrices of dimension 690--713. The correlated window remains bands 12--31. Increasing RKMAX did not increase the 43 stored eigenpairs in the original test. A sum over available complementary bands need not be converged merely because the plane-wave basis is larger.

For a selected band $n$ and a complementary band $a$, the generalized-eigenvector response contains
\begin{equation}
 R_{an}=\frac{\langle a|H'-\epsilon_nO'|n\rangle}{\epsilon_n-\epsilon_a}.
\end{equation}
$H'$ and $O'$ are moving-basis Hamiltonian and overlap derivatives and $\epsilon_n$ are central KS eigenvalues. This diagnostic changes the available complement, holding RKMAX, charge, self-energy, projector window, and normalization input fixed. The implementation includes stored bands both below and above the window; no missing lower-band loop was found.

\begin{center}
\begin{tabular}{rrr}\toprule
Stored bands & Original LAPW0 density term & With installed LAPW0 fixes\\
 & Mean total stress (GPa) & Mean total stress (GPa)\\\midrule
43 & 5.54527 & 5.54761\\
100 & 6.19649 & 6.19884\\
200 & 6.05060 & 6.05295\\
400 & 6.05411 & 6.05645\\
680 & 6.06958 & 6.07193\\\bottomrule
\end{tabular}
\end{center}
The self-energy and correlated window were unchanged. All runs give the same printed SUM and VSUM; their chemical potentials differ by less than $3\times10^{-14}$ eV, and their written SOlapm files have identical hashes. Thus the observed stress change is a response-complement effect, not a changed self-consistent electronic state. The 680-band run still omits the last 10--33 basis eigenstates; it is a near-complete-complement check, not a proof of exact completeness.

The previously verified, charge-relaxed frozen-self-energy finite difference remains $2.93315$ GPa. Its comparison with the original 43-band analytic value is about $2.61$ GPa; with the new density assembly and 680 bands the difference is about $3.14$ GPa. More stored bands do \emph{not} repair the mismatch. Beyond 200 bands the variation is about $0.02$ GPa, much smaller than the remaining discrepancy.

\section{Installed files, reproducibility, and next checks}
The installed executables are \nolinkurl{/Users/haule/wien2k_24/lapw0} and \nolinkurl{/Users/haule/wien2k_24/lcore}. A separate temporary MPI build was tested but no new MPI executable was installed. No dmft2 source or executable was changed during this pass.

The changed source files are:
\begin{itemize}
\item LAPW0: \texttt{modules.F}, \texttt{lapw0.F}, \texttt{xcpot1.f}, \texttt{stress\_gga.F}, \texttt{pwxad5.f}, \texttt{finl\_elect\_str.F}, \texttt{Makefile}, and new \texttt{gga\_metric\_tensor.f}.
\item lcore: \texttt{core\_corr\_stress.f}.
\end{itemize}
Exact backups, source hashes, test inputs/manifests, fixture outputs, and structured result files are in \nolinkurl{src/doc/stress/remaining_terms_audit_20261107/}. The portable supplement in \nolinkurl{wien2k_patches/wien2k24_stress_20261102/supplements/remaining_terms_20261107/} verifies known source hashes, preserves site compiler settings, and makes backups before applying. It does not build or install automatically. It supplements the existing bundle; it is not a claim that the older cumulative installer includes every later G1/G2 or fixed-band-count change.

The remaining work is specific:
\begin{enumerate}
\item Resolve the GGA whole-functional scaling residual by comparing the explicit first derivatives of the energy density with the Euler-potential form, including muffin-tin boundary consistency. Avoid changing the validated pointwise metric tensor to fit material numbers.
\item For SrVO$_3$ LDA, the tested core and electrostatic errors are far too small. Audit the complete one-particle plus normalized-projector variation with matched $H'$, $O'$, and projector derivatives, now using a substantially converged band complement.
\item Isolate the physical strain response of radial functions, energy parameters, and projector radial integrals as a \emph{combined} functional contribution. The previous frozen-radial matrix identities alone do not validate this response. No radial correction has been installed by this audit.
\item Keep numerical reference checks (stationarity, projector convergence, charge, band window and basis membership) separate from matrix-kernel tests. A small combined pressure error can result from cancellation of two incorrect pieces.
\end{enumerate}
The principal SrVO$_3$ eDMFT stress discrepancy remains unresolved. The new tests narrow the search and correct independently demonstrated defects; they do not establish agreement of the complete analytic and numerical stress.
\end{document}
